\documentclass[10pt]{article}
\usepackage[margin=20mm]{geometry}
\usepackage{amsmath,amssymb,mathtools,fontspec,unicode-math,xurl,hyperref}
\IfFontExistsTF{DejaVu Serif}{\setmainfont{DejaVu Serif}}{\setmainfont{Latin Modern Roman}}
\setmathfont{latinmodern-math.otf}
\setlength{\parindent}{0pt}\setlength{\parskip}{6pt}\sloppy
\title{Scientific initial review of
Counting negative eigenvalues for the magnetic Pauli operator}
\author{M25-0144 Initial scientific review}
\date{}
\begin{document}
\maketitle
Abstract: The acquired author manuscript contains a demonstrably invalid trial-space step in its second proof of the principal eigenvalue bound, together with an incorrect illustrative radial-field calculation. The former is classified Major because it directly supplies the central bound along that proof route. A complete local correction is available, and the independent index-theoretic proof is not refuted. The radial calculation is a Minor issue. Explicit checks support the strict spectral thresholds, the boundary-index bookkeeping, the radial equality case and the stated semiclassical deductions. The external scattering trace identity remains an identified, unreproved dependency.\par
\section*{Source, scope and principal statement}
This report concerns S. Fournais, R. L. Frank, M. Goffeng, A. Kachmar and M. Sundqvist, Counting negative eigenvalues for the magnetic Pauli operator, arXiv:2307.16079v1, 29 July 2023, 29 pages [paper]. The assigned journal identifier is DOI `10.1215/00127094-2024-0029`. The final journal text was not acquired; the findings below concern this exact author version. Page numbers refer to its PDF, whose printed pagination agrees.\par
For a smooth bounded connected planar domain with \(d+1\) boundary components, the paper studies the scalar component \(H_{A,g}=(-i\nabla-A)^2-B\) with the specified inward Robin condition. Write \(\Phi_j=(2\pi)^{-1}\int_{\Gamma_j}A_\tau\,\mathrm ds\) and \(\Phi_{g,j}=(2\pi)^{-1}\int_{\Gamma_j}g\,\mathrm ds\), with the boundary orientations fixed in the paper. Theorem 1.1, p. 5, states \[N(H_{A,g},0)\ \geq\ -d+\sum_{j=0}^{d}
          \lceil\Phi_j-\Phi_{g,j}\rceil ,\] where \(N(\,\cdot\,,0)\) counts strictly negative eigenvalues. Equality is asserted on a disk for \(g=0\) and a nonnegative radial field that is nonincreasing in the radius.\par
All 29 pages were read as text. The definitions, quadratic-form identity, radial argument, boundary spectra, adjoint boundary condition, relative-index calculation, conformal reduction and Hardy-space argument received claim-focused checks. Selected source pages were also inspected as rendered images. This is an initial review, not an independent acceptance or a human review.\par
\section*{Established problems}
Problem 1 (Major: the unweighted extension does not give a negative trial space).  In Section 7.2, p. 26, immediately before and in the proof of (7.2), the holomorphic extension \(E(e^{ims})=z^m\) is used to assert \[q_{A,g}(Ev)=\langle P_+D_VP_+v,v\rangle<0,
 \qquad V=g-A_\tau,\] for a vector in the negative spectral subspace of \(P_+D_VP_+\). The ensuing proposed space is the image of that subspace under \(E\). This assertion fails for admissible data.\par
Decisive calculation. Take the unit disk, \(B=\beta>8\), \(g=0\), and \[\phi(z)=\frac{\beta}{4}(|z|^2-1),\qquad
 A=\frac{\beta}{2}(-x_2,x_1).\] These are smooth data in the hypotheses, \(\phi=0\) on the boundary, \(\operatorname{curl}A=\beta\), and \(V=-\beta/2\). The boundary vector \(v=1\) is a strictly negative eigenvector, of eigenvalue \(-\beta/2\), and \(Ev=1\). Direct integration gives \[\langle D_V1,1\rangle=-\pi\beta,\qquad
 q_{A,0}(1)=\int_{D(0,1)}(|A|^2-\beta)\,\mathrm dx
             =\frac{\pi\beta^2}{8}-\pi\beta .\] For example, at \(\beta=16\), the asserted boundary value is \(-16\pi\) while the actual bulk value is \(16\pi>0\). Hence the proposed trial space contains a positive-form vector and is not negative definite. No numerical approximation or limiting regularity issue is involved.\par
Dependency and repair. Proposition 3.1, p. 11, and (5.2), p. 15, have the nonnegative bulk term \[4\|e^{-\phi}\partial_{\bar z}(e^\phi u)\|_{L^2(\Omega)}^2 .\] Ordinary holomorphicity of \(u=Ev\) does not annihilate this term. The required extension is \[E_\phi v=e^{-\phi}Ev .\] For every vector in the finite negative boundary spectral subspace this is an \(H^1\) function, its trace is \(v\), and \(e^\phi E_\phi v=Ev\) is holomorphic. The form identity therefore yields \[q_{A,g}(E_\phi v)=\langle D_Vv,v\rangle<0 .\] Multiplication by the smooth positive factor preserves injectivity. Glazman’s lemma now applies and proves (7.2) with the corrected space.\par
Counterevidence and severity boundary. Remark 3.3, p. 12, already uses \(e^{-\phi}z^m\); Section 5.1 and (5.3), p. 16, also explicitly contain the correct weight. Thus the missing ingredient is present elsewhere in the manuscript. The demonstrated defect is in a step directly establishing the principal bound by the advertised second proof, which warrants Major under the review contract’s core-derivation criterion. It does not show that Theorem 1.1, (7.2), or Lemma 7.3 is false. The local repair above is complete for this step, and the independent APS route remains available.\par
Problem 2 (Minor: the radial field in Remark 2.2 does not give the displayed value).  On p. 10, Remark 2.2 takes \(R=1\), \(m=-1\), and \(B(r)=(r-\tfrac12)^2+\delta\), and assigns the potential at \(r=1\) the value \(1-4/(1+12\delta)^2\), claimed negative for \(0<\delta<1/12\). The actual substitution gives a positive value.\par
Exact substitution. Using the definition \(a(r)=r^{-1}\int_0^r B(\rho)\rho\,\mathrm d\rho\), \[a(1)=\frac1{24}+\frac{\delta}{2},\qquad
 \left(-1-a(1)\right)^2-B(1)
       =\frac{481}{576}+\frac{\delta}{24}+\frac{\delta^2}{4}>0 .\] For \(\delta=1/24\), the actual value is \(643/768\); the printed expression is \(-7/9\). Exact symbolic arithmetic independently confirms these fractions.\par
Repair and impact. Introducing the omitted scale \(B_\beta(r)=\beta[(r-\tfrac12)^2+\delta]\) gives instead \[V_{-1}(1)=1-\frac{\beta}{6}
             +\frac{\beta^2(1+12\delta)^2}{576}.\] Its minimum occurs at \(\beta=48/(1+12\delta)^2\) and equals precisely the printed expression. Alternatively, the smooth Cartesian radial field \(B(r)=12r^4\) gives \(a(1)=2\) and \(V_{-1}(1)=9-12=-3\). Either corrected illustration shows failure of pointwise potential positivity without monotonicity. A negative potential at one point does not itself establish a negative eigenvalue; the original remark also distinguishes these questions. This is an erroneous local illustration, not a counterexample to the radial equality theorem.\par
\section*{Checks supporting the principal argument}
\subsection*{Boundary spectra, orientation and the APS route}
For \(D_V=-i\partial_s+V\) on a component of length \(L\), solving the periodic first-order equation gives \(\mu_m=2\pi(m+\Phi_V)/L\), \(m\in\mathbb Z\). The normal-phase conjugation in Proposition 4.4 has the curvature shift \(V-\kappa\); integration by parts consequently produces the adjoint condition \(\gamma_0u\in(E_{V-\kappa})^\perp\) in Proposition 6.1, p. 19. Strict negativity in \(E_V\) and nonnegativity in its complement are consistent at a zero eigenvalue.\par
The relative-cutoff count in Proposition 6.4, p. 21, contributes \(\lceil-\Phi_{V_j}\rceil\), including at integer flux. The baseline disk index is zero, and the formula \[\frac{1}{4\pi}\int_\Gamma\kappa\,\mathrm ds-\frac{1+d}{2}=-d\] uses \(2\pi\) on the outer component and \(-2\pi\) on each inner component (pp. 22–23). A direct annulus Fourier check gives zero kernel and a one-dimensional cokernel, represented by \(1/\bar z\), hence index \(-1\). These checks support the signs and topology term in the main estimate. For a fixed elliptic boundary space, a smooth change in the bulk magnetic potential is a compact zeroth-order perturbation from \(H^1\) to \(L^2\); this also checks index invariance without assuming that different holonomies are gauge equivalent.\par
There is a local sign slip in Remark 4.2(2), p. 13: on the outer boundary the mean of \(g-A_\tau-\kappa\) is \(\Phi_g-\Phi-1\), not \(\Phi_g-\Phi+1\). The actual shifted eigenvalue computation on p. 15 and the subsequent disk and general-domain calculations use the correct sign. It is not promoted here to an additional mathematical Problem. The general Fredholm and APS theorems invoked by the paper have not been reproved in this review.\par
\subsection*{The radial equality and threshold cases}
For \(m=-k<0\), the radial potential can be written \[\frac{k^2}{r^2}+a^2+\frac{2(k-1)a}{r}
                     +\left(\frac{2a}{r}-B\right),\qquad
 \frac{2a}{r}-B(r)
   =\frac{2}{r^2}\int_0^r[B(\rho)-B(r)]\rho\,\mathrm d\rho .\] It is nonnegative for a nonnegative nonincreasing field. For \(m\geq0\) and the scaled field \(\beta B\), put \(f=r^m e^{-\beta\phi}\), where \(\phi'=a\). The ground-state identity is \[q_m^\beta[u]=
 \int_0^R r f^2\left|(u/f)'\right|^2\,\mathrm dr
       +(m-\beta\Phi)|u(R)|^2 .\] It follows first on the regular core and then on the form domain. If \(m\geq\beta\Phi\), the form is nonnegative. If \(m<\beta\Phi\), \(u=f\) gives a negative direction, whereas the codimension-one trace-zero subspace is nonnegative; there is exactly one negative eigenvalue in that mode. Thus the total is \(\lceil\beta\Phi\rceil\). At an integer flux the threshold mode is zero and is excluded. The zero-field case is nonnegative and is handled separately, without division by \(\Phi=0\). These identities check the conclusion independently of an informal eigenvalue-crossing argument.\par
\subsection*{Semiclassical consequences}
For the magnetic Neumann operator \(\mathcal L_h\) in Theorem 1.4, p. 6, the identity \(\mathcal L_h-h=h^2H_{A/h,0}\) supplies the lower estimate \[N(\mathcal L_h,h)\geq |\Omega|/(2\pi h)+O(1).\] Combined with the cited semiclassical upper estimate, this gives the stated leading term. Corollary 1.6 follows from \[\sum_{e_j<h}e_j=hN(\mathcal L_h,h)-\sum_{e_j<h}(h-e_j).\] The quoted Riesz sum is of order \(h^{1/2}\). On a disk the exact ceiling count makes the error in \(hN\) only \(O(h)\), giving the stated negative \(h^{1/2}\) correction. The external Riesz asymptotic and upper bound were treated as dependencies, not independently reproduced.\par
\section*{Hardy transfer: a domain completion and an external dependency}
The conformal change in Lemma 7.1, pp. 25–26, preserves the magnetic, bulk-field and Robin forms when \(\widetilde A=(DF)^T(A\circ F)\), \(\widetilde B=|F'|^2B\circ F\), and \(\widetilde g=|f'|g\circ f\). The required conclusion is invariance of the negative index under an invertible \(H^1\) map, rather than equality of \(L^2\) norms.\par
On p. 27 the Cayley correspondence explicitly assumes both \(\psi\in L^2(S^1)\) and \(\Psi(t)=\psi((i-t)/(i+t))\in L^2(\mathbb R)\). Thus \(\psi=1\), whose composition is not in \(L^2(\mathbb R)\), does not refute the printed conditional statement. It does reveal why a form-domain argument is needed to transfer arbitrary negative subspaces. One can supply it explicitly. Set \[H_N=\sum_{k=1}^{N}\frac1k,\qquad
 p_N(z)=1-\frac1{H_N}\sum_{k=1}^{N}\frac{(-z)^k}{k}.\] Then \(p_N(-1)=0\), and, with normalized circle measure, \[\|p_N-1\|_2^2=\frac{\sum_{k=1}^{N}k^{-2}}{H_N^2}\longrightarrow0,
 \qquad |p_N-1|_{\dot H^{1/2}}^2=\frac1{H_N}\longrightarrow0.\] Multiplication by a fixed analytic polynomial is continuous in \(H^{1/2}\). Analytic polynomials vanishing at \(-1\) therefore form a dense core in the Hardy form domain. Their Cayley compositions decay as \(O(t^{-1})\), belong to the line Hardy space, and preserve the derivative form and the potential form with \[W(t)=\frac{2}{1+t^2}V\!\left(\frac{i-t}{i+t}\right),\qquad
 \int_{S^1}V|\psi|^2\,\mathrm ds=\int_\mathbb RW|\Psi|^2\,\mathrm dt .\] The last potential identity in the source interchanges \(V\) and \(W\); the formula above is the consistent one. Finite-dimensional strict negativity persists under form-norm approximation. In the reverse direction, the standard weighted unitary Cayley map identifies the Hardy \(L^2\) spaces, and the unweighted inverse composition has an additional bounded factor \(1+z\); conformal invariance of analytic Dirichlet energy controls the derivative form. This establishes both inequalities between the negative indices. It is a completion of the domain justification, not a separate counterexample to Lemma 7.3.\par
For constant \(V=-a\), the circle Hardy eigenvalues are \(m-a\), \(m\geq0\), so the strict negative count is \(\max\{\lceil a\rceil,0\}\). No extra unit is present at integer \(a\). The scattering identity quoted from Kaup–Matsuno in the proof of Lemma 7.3 has a nonnegative integral correction; if valid with its required hypotheses, it implies the claimed ceiling lower bound. Its full Equation (113) was not independently accessible and verified. Selected primary-source scattering definitions and decay assumptions were checked, but this is not a proof of that identity. The Cayley potential decays as \(O(t^{-2})\); compatibility with a weighted space alone does not settle every scattering hypothesis.\par
\section*{Disposition and verification boundary}
Item | Initial-review status\par
Supported assessment | Major: invalid central step of the second proof; one additional Minor calculation.\par
Repair status | Magnetic-weight correction proved locally; radial illustration corrected.\par
Main theorem | No counterexample established; independent APS route survives these findings.\par
Author adoption | Unknown; no author response requested.\par
Remaining work | Check the publisher version and the precise external trace-formula hypotheses; independent adjudication remains pending.\par
Scientific decision | Pending; human reviews: none.\par
Input hashing and text extraction were not counted as scientific proof checks. No author code was executed. No claim of exhaustive reproof of the cited elliptic, index or scattering literature is made.\par
1 S. Fournais, R. L. Frank, M. Goffeng, A. Kachmar and M. Sundqvist, Counting negative eigenvalues for the magnetic Pauli operator, arXiv:2307.16079v1 (29 July 2023), 29 pp. https://arxiv.org/abs/2307.16079v1.\par
\end{document}
